%!TEX root = ../template.tex
%%%%%%%%%%%%%%%%%%%%%%%%%%%%%%%%%%%%%%%%%%%%%%%%%%%%%%%%%%%%%%%%%%%%
%% app-internals.tex
%% NOVA thesis document file
%%
%% Appendix with technical documentation and internals
%%%%%%%%%%%%%%%%%%%%%%%%%%%%%%%%%%%%%%%%%%%%%%%%%%%%%%%%%%%%%%%%%%%%

\chapter{NOVAthesis Internals and Tutorials}
\label{app:internals}

This appendix aggregates technical notes, tutorials, and reference documentation regarding the inner workings of the \gls{novathesis} template, its style files, and configuration mechanisms.

\section{Project Architecture}
\label{sec:project_architecture}

The \novathesis\ template is organized into several directories, each with a specific purpose. Understanding this structure is crucial for any contribution or advanced customization.

\begin{description}
  \item[\texttt{novathesis.cls}] The core \LaTeX\ class file. It handles the loading of packages, defines the document structure, and provides the API for school customization.
  \item[\texttt{0-Config/}] Contains the user-level configuration files.
    \begin{itemize}
      \item \texttt{0\_memoir.tex}: Settings for the \texttt{memoir} class.
      \item \texttt{1\_novathesis.tex}: Main template options (school selection, language, etc.).
      \item \texttt{2\_biblatex.tex}: Bibliography settings.
      \item \texttt{3\_cover.tex}: Manual cover overrides.
      \item \texttt{4\_files.tex}: File inclusion list.
    \end{itemize}
  \item[\texttt{NOVAthesisFiles/}] The "engine" room of the template.
    \begin{itemize}
      \item \texttt{Schools/}: School-specific configuration files (\texttt{.ldf}) and logos.
      \item \texttt{Strings/}: Localization files for different languages.
      \item \texttt{ChapStyles/}: Custom chapter styles.
      \item \texttt{FontStyles/}: Custom font configurations.
      \item[\texttt{StyFiles/}] Internal helper packages (e.g., \texttt{options-ext}, \texttt{nt-covers}, \texttt{memory2}).
      \item \texttt{Images/}: Common images and branding.
    \end{itemize}
\end{description}


\section{Tutorial: Creating a New School Defaults File}
\label{app:school_defaults}

This section provides a detailed guide on how to create a new \texttt{*-defaults.ldf} file to support a new university or school. These files customize the layout, strings, logos, and covers specific to each institution.

The template uses a two-tier naming scheme: \texttt{University / School}. This is reflected in the directory structure under \texttt{NOVAthesisFiles/Schools/}.

\subsection{File Location and Naming}
\label{sec:defaults_location}

All school-specific configuration files should be located in:
\begin{center}
  \texttt{NOVAthesisFiles/Schools/[university]/[school]/}
\end{center}
The filename should follow the convention:
\begin{center}
  \texttt{[university]-[school]-defaults.ldf}
\end{center}

\subsection{Step-by-Step Guide}
\label{sub:adding_school_steps}

\begin{enumerate}
  \item \textbf{Identify the closest match}: Browse the existing directories in \texttt{NOVAthesisFiles/Schools/} to find a school with a similar layout.
  \item \textbf{Create the directory structure}: Create a new directory for your university and school.
  \item \textbf{Create the defaults files}: You may need up to two \texttt{.ldf} files:
    \begin{itemize}
      \item \texttt{.../myuniv/myuniv-defaults.ldf}: For rules common to all schools within the university.
      \item \texttt{.../myuniv/myschool/myuniv-myschool-defaults.ldf}: For school-specific customizations. 
    \end{itemize}
    \textbf{Best Practice}: If rules are consistent across the university, the university-level file should contain the bulk of the logic. If they vary significantly, the school-specific file should be the primary configuration.
  \item \textbf{Add logos}: Create \texttt{Images/} subdirectories at the appropriate level and place logos (preferably in PDF vectorial format).
\end{enumerate}

\subsection{File Structure}
\label{sec:defaults_structure}

A typical \texttt{defaults.ldf} file is divided into several sections.

\subsubsection{Metadata and Global Properties}
At the beginning, define custom options and set template-wide behaviors.

\begin{ltx}
\options{/novathesis,
  myuniv/myschool/custom_option/.new value = {},
}

\ntsetup{print/frontpage=true}
\end{ltx}

\subsubsection{Strings (University, School, and Logos)}
Define the names in multiple languages to support automatic switching. Use the \latexinline{\university} and \latexinline{\school} commands to define localized names.

\begin{ltx}
% University name
\university(pt):={Universidade de Exemplo}
\university(en):={Example University}

% School name
\school(pt):={Faculdade de Testes}
\school(en):={Faculty of Testing}

% Logos
\school(logo,pt):={myschool-logo-pt}
\school(logo,en):={myschool-logo-en}
\school(logo):={\theschool(logo,\@LANG@COVER)}
\end{ltx}

\subsubsection{Colors}
Define school-specific colors using the \texttt{xcolor} package.

\begin{ltx}
\definecolor{schoolblue}{RGB}{0, 50, 100}
\end{ltx}

\subsubsection{Layout and Margins}
Customize margins for normal and cover pages. Margins are defined for different media (\texttt{screen}, \texttt{paper}) and specific pages (\texttt{cover}).

\begin{ltx}
% Normal pages
\margin(screen,top):={3.5cm}
\margin(screen,bottom):={2.5cm}

% Cover pages (1-1 is the front cover, 2-1 is the front page)
\margin(cover,top):={8.6cm}
\margin(cover,bottom):={1.5cm}
\end{ltx}

\subsubsection{Spine Setup}
Configure the appearance of the book spine using \texttt{\string\SpineSetup}.

\begin{ltx}
\SpineSetup{%
  bg/image                    = {myschool-spine-bg},
  box/logo/bg/image           = {\theschool(logo)},
  box/title/fg/color          = {white},
}
\end{ltx}

\subsubsection{Cover Backgrounds}
Set background images or colors for the covers.

\begin{ltx}
\thesiscover(1-1,image):={myschool-cover-front}
\thesiscover(N-2,image):={myschool-cover-back}
\end{ltx}

\section{The \texttt{\textbackslash ntaddtocover} Command}
\label{sec:ntaddtocover_ref}

The primary tool for placing content on the covers: \texttt{\string\ntaddtocover[options]\{covers\}\{code\}}.

\subsection{Command Syntax}
\begin{center}
  \latexinline{\ntaddtocover[<attributes>]{<covers>}{<code>}}
\end{center}

\begin{description}
  \item[\latexinline{attributes}] (Optional) A comma-separated list of key-value pairs that control the layout, alignment, and conditional printing of the element.
  \item[\latexinline{covers}] A comma-separated list of cover identifiers where the code should be applied. Valid identifiers include:
    \begin{itemize}
      \item \texttt{1-1}: Front cover (outer).
      \item \texttt{1-2}: Front cover (inner).
      \item \texttt{2-1}: Front page (inner front cover).
      \item \texttt{2-2}: Back of front page.
      \item \texttt{N-1}: Inside back cover.
      \item \texttt{N-2}: Back cover (outer).
      \item \texttt{spine}: The book spine.
    \end{itemize}
  \item[\latexinline{code}] The actual \LaTeX\ code to be rendered on the specified cover(s).
\end{description}

\subsection{Available Attributes}
\label{sec:attributes}

The behavior of the element can be fine-tuned using the following attributes:

\begin{description}
  \item[\texttt{vspace}] Vertical spacing from the previous element. Supports lengths (\texttt{1cm}) or integers (relative \latexinline{\vfill} weight). Degree-specific variants like \texttt{vspace/phd} are supported.
  \item[\texttt{hspace}] Horizontal spacing.
  \item[\texttt{height} / \texttt{width}] Box dimensions. Special variants like \texttt{height/phd} can be used.
  \item[\texttt{minheight}] Minimum height for the element's container box.
  \item[\texttt{halign} / \texttt{valign}] Alignment (\texttt{l, c, r, j} and \texttt{t, c, b}).
  \item[\texttt{status} / \texttt{degree}] Conditional filters (e.g., \texttt{status=\{working, provisional\}}).
  \item[\texttt{tikz}] If \texttt{true}, code is executed in a page-anchored TikZ environment.
  \item[\texttt{color} / \texttt{bgcolor}] Text and background colors for the container box.
  \item[\texttt{newpar}] Boolean (default \texttt{true}) indicating if a new paragraph should start after rendering.
\end{description}

\subsection{Resetting Covers}
\label{sec:resetting}

If you need to clear all previously defined elements for the covers, you can use the following commands:

\begin{description}
  \item[\latexinline{\ntresetcovers}] Clears all elements from all cover pages.
  \item[\latexinline{\ntresetonecover{<id>}}] Clears elements from a specific page (e.g., \texttt{1-1}).
\end{description}

\subsection{Examples}
\label{sec:examples}

\subsubsection{Adding a logo with TikZ}
This example uses \latexinline{\TikZ} to place a logo at a fixed position:

\begin{ltx}
\ntaddtocover[tikz]{1-1}{%
  \node[anchor=center] at ($(current page.north west) + (20mm,-20mm)$)
    {\includegraphics[width=21mm]{\theuniversity(logo)}};
}
\end{ltx}

\subsubsection{Standard text with spacing and alignment}
This example shows how to add the dissertation title with specific vertical spacing:

\begin{ltx}
\ntaddtocover[vspace=66mm,height=4cm,valign=c]{1-1}{%
  \fontsize{17}{20}\selectfont%
  \textbf{\thedoctitle(\@LANG@COVER,main)}%
}
\end{ltx}

\subsubsection{Conditional elements}
To print a "Draft" timestamp only in working versions:

\begin{ltx}
\ntaddtocover[vspace=2,status=working]{1-1,2-1}{%
  \fontsize{9}{9}\selectfont%
  \emph{\textbf{Draft: \today}}%
}
\end{ltx}


\section{Reference: Metadata and Layout Commands}
\label{sec:metadata_ref}

The template uses a powerful metadata system based on the \texttt{MemStore} package. Most strings can be defined with multiple indices (e.g., degree, gender, language).

\subsection{Core Metadata Commands}
\begin{itemize}
  \item \texttt{\textbackslash school(lang) := \{...\}}: Defines the school name.
  \item \texttt{\textbackslash university(lang) := \{...\}}: Defines the university name.
  \item \texttt{\textbackslash school(logo, lang) := \{filename\}}: Specifies the logo filename.
  \item \texttt{\textbackslash degreestring(type, gender, lang) := \{...\}}: Defines the name of the degree (e.g., \texttt{msc}, \texttt{phd}).
  \item \texttt{\textbackslash majorfield(lang) := \{...\}}: The scientific field or major.
  \item \texttt{\textbackslash specialization(lang) := \{...\}}: Optional specialization.
\end{itemize}

\subsection{Document Descriptions}
\begin{itemize}
  \item \texttt{\textbackslash dissertationstring(type, lang) := \{...\}}: The purpose text on the cover.
  \item \texttt{\textbackslash dissertationplan(type, lang) := \{...\}}: Text used for plans/intermediate reports.
  \item \texttt{\textbackslash copyrighttextstring(lang) := \{...\}}: Legal text for the copyright page.
\end{itemize}

\subsection{Person Titles (Advisers and Committee)}
\begin{itemize}
  \item \texttt{\textbackslash adviserstring(type, num, gender, lang) := \{...\}}:
    \texttt{type} is \texttt{a} (adviser), \texttt{c} (co-adviser), or \texttt{t} (tutor).
  \item \texttt{\textbackslash committeestring(type, num, gender, lang) := \{...\}}: For jury members (e.g., \texttt{c} for chair, \texttt{m} for main examiner).
  \item \texttt{\textbackslash ntprintpersons[options]\{scale\}\{columns\}\{type\}\{order\}}: High-level macro to print lists of persons.
    \begin{itemize}
      \item \texttt{options}: Usually \texttt{name} or \texttt{full}.
      \item \texttt{scale}: Font scale relative to current size.
      \item \texttt{columns}: Number of columns (1 or 2).
      \item \texttt{type}: \texttt{adviser} or \texttt{committee}.
      \item \texttt{order}: Order of roles (e.g., \texttt{c,r,a,ca,m}).
    \end{itemize}
\end{itemize}

\subsection{Layout and Committee Configuration}
...
\begin{itemize}
  \item \texttt{\textbackslash margin(media, side) := \{value\}}: \texttt{media} is \texttt{screen}, \texttt{paper}, or \texttt{cover}.
  \item \texttt{\textbackslash committeeorder() := \{c,m,r,a\}}: Defines the printing order. Keys: \texttt{c} (chair), \texttt{r} (rapporteur), \texttt{a} (adviser), \texttt{ca} (co-adviser), \texttt{m} (member), \texttt{g} (guest).
  \item \texttt{\textbackslash thesiscover(page, attribute) := \{value\}}: Attributes include \texttt{bgcolor}, \texttt{textcolor}, \texttt{image}.
\end{itemize}

\subsection{Utility and Date Commands}
\begin{itemize}
  \item \texttt{\textbackslash ToUppercase\{...\}}: Converts text to uppercase.
  \item \texttt{\textbackslash GetDocDate(type, part)}: Retrieves document dates. \texttt{type} can be \texttt{submission} or \texttt{exam}; \texttt{part} can be \texttt{year}, \texttt{month}, or \texttt{day}.
  \item \texttt{\textbackslash PrintDateISO\{year-month-day\}}: Formats a date according to the template's style.
\end{itemize}


\section{Debugging and Batch Definitions}
\label{sec:debugging_and_loops}

\subsection{Millimeter Grid (Debug Mode)}
You can enable a millimeter grid that covers the entire printing area of the covers and the spine by adding \texttt{debug=cover} to your \latexinline{\ntsetup}.
\begin{ltx}
\ntsetup{debug=cover}
\end{ltx}
This renders a semi-transparent grid with 1mm to 5cm steps to help measure distances.

\subsection{Using a Reference Background Image}
To replicate an existing cover design:
\begin{enumerate}
  \item \textbf{Load the reference image}: Set the background image for the cover page.
  \begin{ltx}
  \thesiscover(1-1,image):={reference-cover}
  \end{ltx}
  \item \textbf{Align your elements}: Use \latexinline{\ntaddtocover} to position text and logos over the reference.
  \item \textbf{Finalize}: Remove the reference image and turn off the debug flag.
\end{enumerate}

\subsection{Batch Definitions with Loops}
To avoid repetitive code for multiple languages or variants, use the internal \latexinline{\ntfor} loop.
\begin{ltx}
\@for\mylang:={\@LANG@LOADED}\do{%
  \committeetitlestring(\mylang):={\ToUppercase{\thecommitteetitlestring(\mylang)}}%
}
\end{ltx}


\section{Localization and Strings}
\label{sec:localization}

The \novathesis\ template is fully multilingual and uses a centralized localization system in \texttt{NOVAthesisFiles/Strings/}.

\subsection{The Localization Mechanism}
The system relies on three interconnected components:
\begin{description}
  \item[\texttt{strings-declare.tex}] Acts as a registry, defining metadata keys that need translation.
  \item[\texttt{strings-<lang>.ldf}] Contains translations for a specific language code.
  \item[\texttt{Babel} Integration] Standard \LaTeX\ labels are managed by \texttt{Babel} but can be overridden using \texttt{\string\ntlangsetup} in the \texttt{.ldf} file.
\end{description}

\subsection{Overriding Babel Labels}
\begin{ltx}
\ntlangsetup*{pt/contents=Índice}
\end{ltx}

\subsection{Adding a New Language}
To add support for a new language code (\texttt{zz}):
\begin{enumerate}
  \item Create \texttt{NOVAthesisFiles/Strings/strings-zz.ldf}.
  \item Input the declarations: \texttt{\string\input\{\string\DIRSTRINGS/strings-declare.tex\}}.
  \item Copy and translate from \texttt{strings-en.ldf}, replacing indices.
\end{enumerate}


\section{Tips \& Tricks}
\label{sec:tips}

\begin{itemize}
  \item \textbf{Expansion}: Use \latexinline{\expanded{...}} when passing memory data (like \latexinline{\theschool(...)}) to commands like \latexinline{\includegraphics}.
  \item \textbf{Hooks}: Inject code at specific points, e.g., \latexinline{\NTAddToHook{cover/1-1/text/pre}{\sffamily}}.
  \item \textbf{Conditionals}: Use \latexinline{\ifphddoc}, \latexinline{\ifmscdoc}, or \latexinline{\ifdocfinal} to tailor the configuration.
  \item \textbf{Utility Commands}: Define small helpers like \latexinline{\printdwithvspace} to simplify layout logic.
  \item \textbf{Loading Fonts}: Use \latexinline{\AtEndPreamble} to load school-specific fonts to avoid clashes.
\end{itemize}


\section{Explanation of copyright.tex}
\label{annex:copyright}

The file \texttt{NOVAthesisFiles/StyFiles/copyright.tex} is a small utility used to generate the default copyright page of the thesis.

\subsection{Functionality}
The file contains the necessary commands to format and print the copyright information:
\begin{itemize}
    \item \textbf{Vertical Alignment}: It uses \texttt{\textbackslash null\textbackslash vfill} to push the copyright text to the bottom of the page.
    \item \textbf{Document Title}: It prints the document title using the \texttt{\textbackslash thedoctitle} command.
    \item \textbf{Copyright String}: It invokes \texttt{\textbackslash thecopyrightstr}, which contains the localized legal text.
\end{itemize}

\subsection{Summary}
In summary, \texttt{copyright.tex} provides a standardized way to render the copyright page, maintaining consistency across all documents produced with the \emph{NOVAthesis} template.


\section{Explanation of fix-babel-expl3.tex}
\label{annex:fix-babel-expl3}

The file \texttt{NOVAthesisFiles/StyFiles/fix-babel-expl3.tex} is a modern refactoring of the original localization utility. It uses \LaTeX3 (\texttt{expl3}) to provide a more robust way to manage localized strings (captions) across multiple languages.

\subsection{Core Functionality}
The primary entry point is the \texttt{\textbackslash ntlangsetup} command.
\begin{itemize}
    \item \textbf{Syntax}: \texttt{\textbackslash ntlangsetup*[suffix]\{lang/key=value\}}
    \item \textbf{Property Lists}: Internally uses \texttt{expl3} property lists (\texttt{prop}) named \texttt{g\_nt\_babel\_<lang>\_prop}.
    \item \textbf{Application}: Collected translations are applied at the end of the preamble using the \texttt{begindocument/before} hook and Babel's caption system.
\end{itemize}

\subsection{Summary}
By moving to \texttt{expl3}, \emph{NOVAthesis} gains a cleaner separation between translation definition and application, resulting in a more maintainable localization system.


\section{Explanation of fix-bib.bib}
\label{annex:fix-bib}

The file \texttt{NOVAthesisFiles/StyFiles/fix-bib.bib} contains mandatory bibliography entries for the template, often referenced in the manual and examples.

\subsection{Template Documentation}
The most important entry is \texttt{novathesis-manual}, which provides the official reference for the template. Users are explicitly requested to \texttt{\textbackslash cite\{novathesis-manual\}} in their document.

\subsection{Summary}
The \texttt{fix-bib.bib} file ensures that critical references for the template, \LaTeX, and AI ethics are always available and correctly formatted.


\section{Explanation of fix-hyperref.sty}
\label{annex:fix-hyperref}

The file \texttt{NOVAthesisFiles/StyFiles/fix-hyperref.sty} extends the \texttt{hyperref} package to fix language mapping issues and provide a capitalized version of the \texttt{\textbackslash autoref} command.

\subsection{The \textbackslash Autoref Command}
Standard \latexinline{\autoref} usually prints labels in lowercase. This file implements a robust \texttt{\textbackslash Autoref} command using \LaTeX3 (\texttt{Expl3}) logic that:
\begin{enumerate}
    \item Automatically detects the target type (chapter, section, etc.).
    \item Fetches the localized name.
    \item \textbf{Capitalizes the first letter} using an internal helper.
\end{enumerate}

\subsection{Summary}
In summary, this file ensures that cross-references work reliably across all languages and provide correctly capitalized localized references.


\section{Explanation of glossary-xltabular.sty}
\label{annex:glossary-xltabular}

The file \texttt{NOVAthesisFiles/StyFiles/glossary-xltabular.sty} is a custom extension for the \texttt{glossaries} package, providing new tabular styles using the \texttt{xltabular} package.

\subsection{Key Features}
The package defines sophisticated styles (e.g., \texttt{xltabular}, \texttt{xltabular4col}) for professional, tabular glossaries. These styles integrate with the template's configuration system, allowing users to change the look simply by adjusting a key-value option.

\subsection{Summary}
In summary, \texttt{glossary-xltabular.sty} provides high-quality, multi-page tabular layouts for lists of acronyms and symbols.


\section{Explanation of memory2.sty}
\label{annex:memory2}

The file \texttt{NOVAthesisFiles/StyFiles/memory2.sty} is a legacy utility package used for storing and retrieving structured data within \LaTeX{}.

\subsection{Key Concepts}
It allows the creation of data structures that resemble objects or arrays:
\begin{itemize}
    \item \textbf{\textbackslash newdata}: Defines a new data storage object.
    \item \textbf{\textbackslash the<Name>(key)}: Retrieves a stored value.
\end{itemize}

\subsection{Summary}
While newer parts of the template favor \texttt{memstore.sty}, \texttt{memory2.sty} is still used for various legacy features and administrative checklists.


\section{Explanation of memstore.sty}
\label{annex:memstore}

The file \texttt{NOVAthesisFiles/StyFiles/memstore.sty} is a modern memory storage engine for \emph{NOVAthesis}, built using \LaTeX3 (\texttt{expl3}).

\subsection{Core Commands}
\begin{itemize}
    \item \textbf{\textbackslash MemSet} / \textbf{\textbackslash MemGet}: Stores and retrieves values under keys in named containers.
    \item \textbf{\textbackslash NewMemStore}: Generates custom interfaces like \texttt{\string\SetAuthors}.
    \item \textbf{\textbackslash IfMemVoidTF}: Conditional helper for detecting missing or blank values.
\end{itemize}

\subsection{Summary}
In summary, \texttt{memstore.sty} is the state-of-the-art data engine for \emph{NOVAthesis}, providing underlying infrastructure for managing complex document metadata.


\section{Explanation of novathesis-spine.sty}
\label{annex:novathesis-spine}

The file \texttt{NOVAthesisFiles/StyFiles/novathesis-spine.sty} provides a configurable engine for generating thesis book spines using \texttt{expl3}, \texttt{TikZ}, and \texttt{tcolorbox}.

\subsection{Core Commands}
\begin{itemize}
    \item \textbf{\textbackslash SpineSetup}: Configures spine parameters via key-value interface.
    \item \textbf{\textbackslash PrintBookSpine}: Typesets the horizontal spine strip, automatically adjusting width and thickness.
\end{itemize}

\subsection{Summary}
In summary, \texttt{novathesis-spine.sty} is a professional-grade tool that automates the generation of polished book spines.


\section{Explanation of options-ext.sty}
\label{annex:options-ext}

The file \texttt{NOVAthesisFiles/StyFiles/options-ext.sty} provides powerful extensions to the \texttt{options} package, crucial for the template's complex key-value management system.

\subsection{Key Handlers}
This package adds handlers for choice iteration, list copying (\texttt{.copy list}), and conditional pushing (\texttt{.newpush ifnew}). These extensions allow the template's configuration files to be highly expressive and declarative.

\subsection{Summary}
In summary, \texttt{options-ext.sty} is the engine that enables \emph{NOVAthesis} to handle its highly modular and customizable architecture.


\section{Explanation of nt-covers.sty}
\label{annex:nt-covers}

The file \texttt{NOVAthesisFiles/StyFiles/nt-covers.sty} contains the core engine for generating and customizing the thesis covers and front pages. It was isolated from the main class file to improve maintainability.

\subsection{Core Functionality}
This package manages:
\begin{itemize}
    \item \textbf{\textbackslash ntaddtocover}: The primary interface for adding elements to covers.
    \item \textbf{Cover Printing}: Macros like \texttt{\textbackslash ntprintcovers} and \texttt{\textbackslash ntprintcoverpair} that orchestrate the rendering of the different cover layers.
    \item \textbf{Dynamic Layout}: The logic for handling margins, backgrounds, and element positioning (including TikZ overlays).
    \item \textbf{Conditional Rendering}: Filtering elements based on document status or degree type.
\end{itemize}

\subsection{Summary}
By centralizing cover-related logic in \texttt{nt-covers.sty}, the template maintains a clean separation between document structure and the complex graphical requirements of institution-specific covers.


\section{Explanation of paralist-compat.tex}
\label{annex:paralist-compat}

The file \texttt{NOVAthesisFiles/StyFiles/paralist-compat.tex} provides custom list environments that emulate the behavior of the popular \texttt{paralist} package using the modern \texttt{enumitem} package.

\subsection{Emulated Environments}
It defines environments like \texttt{compactitem} and \texttt{inparaenum} to ensure compatibility with styles that expect these features, while maintaining the technical advantages of \texttt{enumitem}.

\subsection{Summary}
In summary, \texttt{paralist-compat.tex} acts as a lightweight shim for familiar list environments.
